\section{\qII}
\newcommand{\boxplotfig}{
    \begin{figure}
        
    \end{figure}
}

\subsection{Introduction}
This subquestion investigates empirical data regarding the I - V - IV - iv and other chord progressions, by analyzing the results of an online survey. This yields an indication of the extent of perceived enjoyment upon initial, isolated observation.

\subsection{Results}
% RESULTS ANALYZED FROM 18/09/2026
Each fragment is labeled by a letter. % Refer to the following list to find the corresponding chord progression of each letter:
% pictures of all graphs

% graph with an x-axis of "fragment" and y-axis of "perceived enjoyment", with points denoting the average of each fragment.

\subsection{Analysis (discrete)}
The boxplots show that fragments C and J share the highest distribution of ratings overall. Fragment J has a higher average of approximately $4$, while fragment C has a lower average of approximately $3$.

Fragment C shares the same mean with fragments B, D, E, F, H, and I.

\subsection{Analysis (continuous)}
The results yield an inverted-U shaped relation of the perceived enjoyment as harmonic complexity increases, with only one outlier. 

For fragments A to C, the average rated enjoyment increases, and fragment C reaches the highest distribution of ratings. However, the following fragment D has a more centered distribution than C, signifying the start of a decline. While the following fragment E has a lower distribution of ratings, fragment F moves back to a centered distribution of 3. Then, fragment G has the lowest distribution of ratings, with almost all values being 1.

Fragments H to J, having a similar harmonic complexity to fragment C, also have a similar average perceived enjoyment. The mean of fragments H and I are equal to that of fragment C, though their distributions are lower; fragment J has a higher mean, but equal distribution to that of fragment C.

\subsection{Conclusion (discrete)}
% The discrete conclusion draws a conclusion from results analyzed in this subquestion.

Given that the I - V - IV - iv chord progression has the highest distribution of ratings overall, we can infer that the I - V - IV - iv has a moderately high shared perceived enjoyment upon initial listen. 

It is not uniquely higher than all other chord progressions, though; the I - vi - IV - V chord progression has not only a congruent distribution of ratings, but a higher mean, from which we can infer an even higher shared perceived enjoyment. Also, many other chord progression had an approximately congruent mean as the I - V - IV - iv chord progression, despite its higher distribution.

\subsection{Conclusion (continuous)}
% The continuous conclusion draws a conclusion using data from subquestion 1 as well as this subquestion.

Given that there is an inverted U-shaped relation between the harmonic complexity of chord progressions and the perceived enjoyment, with I - V - IV - iv having the overall highest distribution and mean of perceived enjoyment, we can infer that the I - V - IV - iv sits at the vertex of this shape. Our conclusion from subquestion 1, claiming that the I - V - IV - iv chord progression sits at a "sweet spot" of harmonic complexity, and that — [according to the Wundt Cuvre] — [inverted u shape]